% !TEX TS-program = pdflatex
\documentclass[11pt]{article}
\usepackage[utf8]{inputenc}
\usepackage[T1]{fontenc}
\usepackage{lmodern}
\usepackage{geometry}
\geometry{margin=1in}
\usepackage{microtype}
\usepackage{amsmath, amssymb, amsthm, mathtools}
\usepackage{mathrsfs}
\usepackage{bm}
\usepackage{enumitem}
\usepackage{graphicx}
\usepackage{booktabs}
\usepackage{xcolor}
\usepackage{hyperref}
\usepackage[nameinlink,capitalize]{cleveref}
\usepackage{csquotes}
\usepackage{tikz}
\usepackage{caption}
\usepackage{subcaption}
% Optional for code/algorithms
\usepackage{listings}
\lstset{basicstyle=\ttfamily\small,breaklines=true,columns=fullflexible,frame=single,numberstyle=\tiny,numbers=left}

%----- theorem styles -----
\theoremstyle{plain}
\newtheorem{theorem}{Theorem}[section]
\newtheorem{lemma}[theorem]{Lemma}
\newtheorem{proposition}[theorem]{Proposition}
\newtheorem{corollary}[theorem]{Corollary}
\theoremstyle{definition}
\newtheorem{definition}[theorem]{Definition}
\theoremstyle{remark}
\newtheorem{remark}[theorem]{Remark}

%----- macros -----
\DeclareMathOperator{\vtwo}{v_2}
\newcommand{\Z}{\mathbb{Z}}
\newcommand{\N}{\mathbb{N}}
\newcommand{\Odd}{\mathcal{O}} % odd positive integers
\newcommand{\Todd}{T} % odd macro map
\newcommand{\V}{V}
\newcommand{\W}{W}
\newcommand{\wt}{\widetilde}
\newcommand{\E}{\mathbb{E}}
\newcommand{\1}{\mathbf{1}}

\title{Collatz via Drift on Odd Macro States: A Rigorous Core with Visual Intuition}
\author{Ryan O. Van Gelder \& Luis Morat´o de Dalmases}
\date{\today}

\begin{document}
\maketitle

 
\begin{abstract}
We present a compact, rigorous reduction of the Collatz termination problem to a frequency surplus condition for the $2$-adic valuations of odd macro steps. The analytic core is an exact per-step drift identity for $\V(n)=\log n$ along the odd-state map $\Todd$, together with ambient density facts for $\vtwo(3n+1)$. We then state a conditional theorem: if the time-averaged cumulative frequency of events $\{\vtwo(3n_t+1)\ge k\}$ exceeds their ambient $2$-adic densities by a total surplus greater than a small universal buffer, then every trajectory terminates and no nontrivial odd cycle exists. We complement this with visualization-only figures (parity tree, affine branches on $\hat{\mathbb C}$) to build intuition for the counting quantities actually used in the proof, and with an empirical section reporting frequency histograms versus the ambient benchmark. All prior claims of $2$-adic contractivity are removed; visuals are explicitly labeled as heuristic.
\end{abstract}

\section{Roadmap and Objects}
Let $\Odd=\{n\in\N: n\text{ odd}\}$. Define the \emph{odd macro map}
\begin{equation}
  \Todd(n) \coloneqq \frac{3n+1}{2^{\vtwo(3n+1)}},\qquad n\in\Odd,
\end{equation}
which maps odd integers to odd integers by executing a Collatz ``odd step'' followed by the maximal division by $2$.
Set the Lyapunov candidate $\V(n)\coloneqq\log n$ (natural logarithm), and write the induced trajectory $n_{t+1}=\Todd(n_t)$ with $n_0\in\Odd$.

We also adopt a convenient unit system:
\begin{equation}
  \gamma \coloneqq \frac{\log 3}{\log 2},\qquad \theta_{\max} \coloneqq \frac{\log(4/3)}{\log 2},\qquad \text{so that }\;\gamma+\theta_{\max}=2.
\end{equation}
We measure per-step drift both in natural units and in ``base-2 units,'' where dividing by $\log 2$ simplifies bookkeeping.

\section{Exact One-Step Drift and a Uniform Remainder}
\begin{lemma}[Drift identity]
For $n\in\Odd$,
\begin{equation}\label{eq:drift}
  \V(\Todd(n)) - \V(n) 
  \,=\, \log 3 \, -\, \vtwo(3n+1)\,\log 2 \, +\, \log\!\Bigl(1+\tfrac{1}{3n}\Bigr).
\end{equation}
Equivalently, in base-2 units with $\W\coloneqq \V/\log 2$,
\begin{equation}\label{eq:drift-base2}
  \W(\Todd(n)) - \W(n)
  \,=\, \gamma \, -\, \vtwo(3n+1) \, +\, \theta(n),\qquad \theta(n)\coloneqq \frac{\log(1+1/(3n))}{\log 2}.
\end{equation}
Moreover, $0<\theta(n)\le \theta_{\max}$ with equality only at $n=1$.
\end{lemma}
\begin{proof}
Immediate from $\Todd(n)=(3n+1)/2^{\vtwo(3n+1)}$ and $\log(3n+1)=\log 3+\log n+\log(1+1/(3n))$, followed by the bound $\log(1+x)\le x$ with $x\in(0,1/3]$ giving $\log(1+1/(3n))\le\log(4/3)$.
\end{proof}

\section{Ambient $2$-Adic Densities and a Counting Identity}
For $k\ge1$ and $n\in\Odd$, define the indicators
\begin{equation}
  f_k(n) \coloneqq \1\{\vtwo(3n+1)\ge k\} .
\end{equation}
Then the valuation decomposes as the convergent sum $\vtwo(3n+1)=\sum_{k\ge1} f_k(n)$. The following facts are standard.

\begin{proposition}[Ambient densities and mean]
Fix $k\ge1$. Exactly one odd residue class modulo $2^k$ satisfies $\vtwo(3n+1)\ge k$. Consequently the upper/lower natural densities of $\{n\in\Odd: \vtwo(3n+1)\ge k\}$ inside the odd integers both equal $2^{-(k-1)}$. In particular,
\begin{equation}
  \lim_{X\to\infty} \frac{1}{\#\{n\le X\,:\, n\text{ odd}\}} \sum_{\substack{n\le X\\ n\,\text{odd}}} f_k(n) \,=\, 2^{-(k-1)},\qquad \sum_{k\ge1} 2^{-(k-1)} \,=\, 2.
\end{equation}
\end{proposition}

\section{A Conditional Drift Theorem}
Let $n_{t+1}=\Todd(n_t)$ and define temporal averages
\begin{equation}
  \alpha_k \coloneqq \liminf_{T\to\infty} \frac{1}{T} \sum_{t<T} f_k(n_t) \;\in[0,1].
\end{equation}
Note that by Fatou's lemma for nonnegative series,
\begin{equation}\label{eq:fatou}
  \liminf_{T\to\infty} \frac{1}{T} \sum_{t<T} \sum_{k\ge1} f_k(n_t) \;\ge\; \sum_{k\ge1} \alpha_k.
\end{equation}

\begin{theorem}[Frequency surplus $\Rightarrow$ termination]
Suppose there exists $\eta>0$ such that
\begin{equation}\label{eq:surplus}
  \sum_{k\ge1} \alpha_k \;\ge\; 2+\eta. 
\end{equation}
Then no Collatz orbit in $\Odd$ is infinite, and there are no nontrivial odd cycles. Equivalently, every trajectory enters the trivial cycle $1\mapsto1$.
\end{theorem}
\begin{proof}
Summing \eqref{eq:drift-base2} along the trajectory and dividing by $T$ gives
\begin{align}
  \frac{\W(n_T)-\W(n_0)}{T}
  &= \gamma - \frac{1}{T}\sum_{t<T} \sum_{k\ge1} f_k(n_t) + \frac{1}{T}\sum_{t<T} \theta(n_t) \\
  &\le \gamma - \frac{1}{T}\sum_{t<T} \sum_{k\ge1} f_k(n_t) + \theta_{\max} .
\end{align}
Taking $\limsup$ on the left and $\liminf$ on the right, and using \eqref{eq:fatou} plus $\gamma+\theta_{\max}=2$, we obtain
\begin{equation}
  \limsup_{T\to\infty} \frac{\W(n_T)-\W(n_0)}{T} \;\le\; 2 - \sum_{k\ge1} \alpha_k \;\le\; -\eta.
\end{equation}
Thus for all large $T$, $\W(n_T) \le \W(n_0) - \eta T$, i.e. $\V(n_T) \le \V(n_0) - (\eta\log2)\,T$. Since $\V(n_T)=\log n_T\ge0$ for all $T$, this is impossible unless the trajectory stops growing and enters the cycle at $1$. Hence no infinite orbit and no nontrivial odd cycle can exist under \eqref{eq:surplus}.
\end{proof}

\begin{remark}[Interpretation]
The identity $\vtwo(3n+1)=\sum_k f_k(n)$ converts termination into a purely \emph{counting} surplus problem: the ambient (random-like) benchmark contributes exactly $2$ in base-2 units, while the small remainder $\theta(n)\in(0,\theta_{\max}]$ provides a universal buffer that the surplus must exceed.
\end{remark}

\section{Visualization Section (Heuristic Only)}
\emph{This section is for intuition; no figure is used in proofs.} We illustrate:
\begin{enumerate}[label=\textbf{V\arabic*}.,leftmargin=*]
  \item The parity tree with branches labeled by the affine maps $n\mapsto n/2$ and $n\mapsto (3n+1)/2$.
  \item Snapshots of these affine maps as Möbius transformations on the Riemann sphere $\hat{\mathbb C}$ (for geometric context only).
  \item Residue plots for $n\bmod 2^k$ colored by $k=\vtwo(3n+1)$, visually connecting to the indicators $f_k$.
\end{enumerate}
\begin{figure}[h]
  \centering
  \fbox{\rule{0pt}{1.75in}\rule{0.95\linewidth}{0pt}}\\[-2pt]
  \caption{Parity tree and affine branches (placeholder).}
\end{figure}
\begin{figure}[h]
  \centering
  \fbox{\rule{0pt}{1.75in}\rule{0.95\linewidth}{0pt}}\\[-2pt]
  \caption{Residue-class coloring by $k=\vtwo(3n+1)$ (placeholder).}
\end{figure}

\section{Empirical Checks Against Ambient Benchmarks}
We report simple computations designed as \emph{sanity checks} for the analytic targets above. For each seed $n_0\in\Odd$ and horizon $T$:
\begin{enumerate}[label=\arabic*.]
  \item Generate $n_{t+1}=\Todd(n_t)$ and record $k_t\coloneqq \vtwo(3n_t+1)$.
  \item For each $k\in\{1,\dots,K\}$, accumulate counts of events $\{k_t\ge k\}$, producing empirical frequencies $\hat f_k(T)=\frac{1}{T}\sum_{t<T}\1\{k_t\ge k\}$.
  \item Compare $\hat f_k(T)$ to $2^{-(k-1)}$ via bar plots and report the cumulative surplus $\sum_{k\le K}\bigl(\hat f_k(T)-2^{-(k-1)}\bigr)$.
\end{enumerate}
\noindent Example Python-like pseudocode (for reproducibility; replace with your environment):
\begin{lstlisting}[language=Python]
from math import log

def v2(m):
    c = 0
    while m % 2 == 0:
        m //= 2
        c += 1
    return c

def T(n):
    return (3*n+1) >> v2(3*n+1)

def run(seed=27, Tsteps=2_000_000, K=12):
    n = seed
    counts = [0]* (K+1)  # 1..K used
    for t in range(Tsteps):
        k = v2(3*n+1)
        for j in range(1, min(k, K)+1):
            counts[j] += 1
        n = T(n)
    freqs = [counts[j]/Tsteps for j in range(1, K+1)]
    ambient = [2**(-(j-1)) for j in range(1, K+1)]
    surplus = sum(f - a for f, a in zip(freqs, ambient))
    return freqs, ambient, surplus
\end{lstlisting}
\begin{remark}
We emphasize: empirical checks serve only to illustrate how one measures the frequencies that appear in the theorem. They are \emph{not} evidence of contractivity and do not substitute for the surplus condition \eqref{eq:surplus}.
\end{remark}

\section{Optional: Koopman/Transfer Operator Viewpoint}
Define the Koopman operator $\mathcal{U}$ on bounded functions $\varphi:\Odd\to\mathbb{R}$ by $(\mathcal{U}\varphi)(n)=\varphi(\Todd(n))$. Choosing $\varphi=\log$ and expanding $\log(3n+1)$ yields the drift identity \eqref{eq:drift}. This framing can aid in connecting drift to spectral questions on function spaces; we do not pursue a functional-analytic theory here.

\section*{Checklist for Figures and Claims}
\begin{itemize}
  \item All figures from the visualization section are labeled \emph{heuristic only}. No analytic claim relies on them.
  \item No $2$-adic contractivity statements are made.
  \item The averaged-drift inequality uses the corrected factor $-\eta\,\log 2$ (no extraneous $1/2$).
  \item Frequencies are compared against the ambient benchmark $2^{-(k-1)}$; cumulative surplus is the reported scalar.
\end{itemize}

\section*{Acknowledgments and Notes}
This document merges a rigorous drift reduction with curated visuals/code from an earlier exploratory draft. All cross-topology language (Riemann sphere vs. $2$-adics) is kept purely illustrative.

%----------------------------

\nocite{*}
\bibliographystyle{unsrt}
\bibliography{references}

\end{document}
